\subsection*{Questão 11}

\begin{questionbox}{11}
  A reta $r$ passa pelo ponto $P(1,4,3)$ e é paralela ao vetor $\mathbf{v} = (-3,-2,-3)$. Determine o ponto de interseção da reta $r$ com o plano $xy$.
\end{questionbox}

\newpage

\begin{solutionbox}
  A equação paramétrica da reta é dada por:
  \[
  (x,y,z) = (1,4,3) + t(-3,-2,-3).
  \]

  Logo,
  \[
  x = 1 - 3t,\quad y = 4 - 2t,\quad z = 3 - 3t.
  \]

  O plano $xy$ é definido por $z = 0$. Assim:
  \[
  3 - 3t = 0 \Rightarrow t = 1.
  \]

  Substituindo em $x$ e $y$:
  \[
  x = 1 - 3(1) = -2,\quad y = 4 - 2(1) = 2.
  \]
  \end{solutionbox}

  \begin{solutionbox}
  Portanto, o ponto de interseção é:
  \[
  \boxed{(-2,\,2,\,0)}.
  \]
\end{solutionbox}